% GENERATED FILE. Edit canonical lessons, then run npm run generate:books.
% canonical-source-hash: fnv1a64:0e09fa5bfd315549
% canonical-lessons: KA-C371-tali, KA-C371-pallakki, KA-C371-palegara, KA-C371-mahalu, KA-C371-kalidu

\chapter{The Chieftain Comes Home}
\label{ch:ka-the-chieftain-comes-home}
\hlchaptermodality{371}

\begin{chapteropening}
\textbf{By the end of this chapter:} \emph{I can say a wedding pendant, a palanquin, a chieftain, a mansion and to set foot.}

Complete the last lesson of chapter 371: I can say a wedding pendant, a palanquin, a chieftain, a mansion and to set foot.
\end{chapteropening}

\section[\texorpdfstring{\kn{ತಾಳಿ} (tāḷi)}{ತಾಳಿ (tāḷi)}]{\kn{ತಾಳಿ} (tāḷi) — a wedding pendant}

\label{lesson:KA-C371-tali}



\begin{quote}
Before the new one: say the Kannada for a listener, then the Kannada for notorious.
\end{quote}

\subsection*{You'll want to know: \kn{ತಾಳಿ}}
\textbf{\kn{ತಾಳಿ}} — \emph{tāḷi} — \textquotedblleft{}a wedding pendant\textquotedblright{}.

\begin{quote}
\kn{ಮದುವೆಯಲ್ಲಿ} \kn{ಗಂಡು} \kn{ಹೆಣ್ಣಿಗೆ} \kn{ತಾಳಿ} \kn{ಕಟ್ಟಿದನು}.
\end{quote}

\begin{quote}
\emph{maduveyalli gaṇḍu heṇṇige tāḷi kaṭṭidanu.}
\end{quote}

\begin{quote}
At the wedding the groom tied the wedding pendant on the bride.
\end{quote}

The gold pendant on a cord tied at a Hindu wedding; \kn{ತಾಳಿ} \kn{ಕಟ್ಟು}, to tie it, means to marry.

\begin{cousinweb}
\textbf{\kn{ತಾಳಿ}} (\emph{tāḷi}) is native Dravidian: DEDR 3175 lists Kannada \emph{tāli}, \emph{tāḷi}, \textquotedblleft{}a small round plate of gold worn at the neck as a marriage badge\textquotedblright{}, beside Tamil \emph{tāli}.
\end{cousinweb}

\begin{grammarlens}[title={What you've built}]
One more word for telling what happened.
\end{grammarlens}

\subsection*{Your turn}
\begin{itemize}
\raggedright
  \item \emph{Say it:} \emph{tāḷi}
  \item \emph{Say it:} \emph{tāḷi}, once more
  \item \emph{Say it:} \emph{maduveyalli gaṇḍu heṇṇige tāḷi kaṭṭidanu.} — At the wedding the groom tied the wedding pendant on the bride.
  \item \emph{Recall:} say the Kannada for a listener, then the Kannada for notorious, then say \emph{tāḷi} again
\end{itemize}

\begin{tcolorbox}[breakable,colback=teal!4,colframe=teal!35!black,title={Before you move on}]
What does \kn{ತಾಳಿ} mean? (\textquotedblleft{}A wedding pendant\textquotedblright{}.) Say it once more.
\end{tcolorbox}
\section[\texorpdfstring{\kn{ಪಲ್ಲಕ್ಕಿ} (pallakki)}{ಪಲ್ಲಕ್ಕಿ (pallakki)}]{\kn{ಪಲ್ಲಕ್ಕಿ} (pallakki) — a palanquin}

\label{lesson:KA-C371-pallakki}



\begin{quote}
Before the new one: say the Kannada for notorious, then the Kannada for a wedding pendant.
\end{quote}

\subsection*{You'll want to know: \kn{ಪಲ್ಲಕ್ಕಿ}}
\textbf{\kn{ಪಲ್ಲಕ್ಕಿ}} — \emph{pallakki} — \textquotedblleft{}a palanquin\textquotedblright{}.

\begin{quote}
\kn{ರಾಣಿ} \kn{ಪಲ್ಲಕ್ಕಿಯಲ್ಲಿ} \kn{ದೇವಸ್ಥಾನಕ್ಕೆ} \kn{ಹೋದಳು}.
\end{quote}

\begin{quote}
\emph{rāṇi pallakkiyalli dēvasthānakke hōdaḷu.}
\end{quote}

\begin{quote}
The queen went to the temple in a palanquin.
\end{quote}

Plural: \kn{ಪಲ್ಲಕ್ಕಿಗಳು}, \emph{pallakkigaḷu}.

A covered seat carried on poles by bearers.

\begin{cousinweb}
\textbf{\kn{ಪಲ್ಲಕ್ಕಿ}} (\emph{pallakki}) is usually traced to Sanskrit \emph{palyaṅka}, \emph{paryaṅka}, \textquotedblleft{}a bed, a couch, a litter\textquotedblright{} (CDIAL 7964).
\end{cousinweb}

\begin{grammarlens}[title={What you've built}]
One more word for telling what happened.
\end{grammarlens}

\subsection*{Your turn}
\begin{itemize}
\raggedright
  \item \emph{Say it:} \emph{pallakki}
  \item \emph{Say it:} \emph{pallakki}, once more
  \item \emph{Say it:} \emph{rāṇi pallakkiyalli dēvasthānakke hōdaḷu.} — The queen went to the temple in a palanquin.
  \item \emph{Recall:} say the Kannada for notorious, then the Kannada for a wedding pendant, then say \emph{pallakki} again
\end{itemize}

\begin{tcolorbox}[breakable,colback=teal!4,colframe=teal!35!black,title={Before you move on}]
What does \kn{ಪಲ್ಲಕ್ಕಿ} mean? (\textquotedblleft{}A palanquin\textquotedblright{}.) Say it once more.
\end{tcolorbox}
\section[\texorpdfstring{\kn{ಪಾಳೆಗಾರ} (pāḷegāra)}{ಪಾಳೆಗಾರ (pāḷegāra)}]{\kn{ಪಾಳೆಗಾರ} (pāḷegāra) — a chieftain; a local lord}

\label{lesson:KA-C371-palegara}



\begin{quote}
Before the new one: say the Kannada for a wedding pendant, then the Kannada for a palanquin.
\end{quote}

\subsection*{You'll want to know: \kn{ಪಾಳೆಗಾರ}}
\textbf{\kn{ಪಾಳೆಗಾರ}} — \emph{pāḷegāra} — \textquotedblleft{}a chieftain; a local lord\textquotedblright{}.

\begin{quote}
\kn{ಆ} \kn{ಕೋಟೆಯನ್ನು} \kn{ಒಬ್ಬ} \kn{ಪಾಳೆಗಾರ} \kn{ಕಟ್ಟಿದನು}.
\end{quote}

\begin{quote}
\emph{ā kōṭeyannu obba pāḷegāra kaṭṭidanu.}
\end{quote}

\begin{quote}
A local chieftain built that fort.
\end{quote}

Plural: \kn{ಪಾಳೆಗಾರರು}, \emph{pāḷegāraru}.

A lord of a small territory under a king; the English wrote \textquotedblleft{}poligar\textquotedblright{}.

\begin{cousinweb}
\textbf{\kn{ಪಾಳೆಗಾರ}} (\emph{pāḷegāra}) is native Dravidian: it is built on \emph{pāḷeya}, \textquotedblleft{}an army camp, a fief\textquotedblright{} (DEDR 4117), with the ending -\emph{gāra}, \textquotedblleft{}one who has, one who does\textquotedblright{}.
\end{cousinweb}

\begin{grammarlens}[title={What you've built}]
One more word for telling what happened.
\end{grammarlens}

\subsection*{Your turn}
\begin{itemize}
\raggedright
  \item \emph{Say it:} \emph{pāḷegāra}
  \item \emph{Say it:} \emph{pāḷegāra}, once more
  \item \emph{Say it:} \emph{ā kōṭeyannu obba pāḷegāra kaṭṭidanu.} — A local chieftain built that fort.
  \item \emph{Recall:} say the Kannada for a wedding pendant, then the Kannada for a palanquin, then say \emph{pāḷegāra} again
\end{itemize}

\begin{tcolorbox}[breakable,colback=teal!4,colframe=teal!35!black,title={Before you move on}]
What does \kn{ಪಾಳೆಗಾರ} mean? (\textquotedblleft{}A chieftain; a local lord\textquotedblright{}.) Say it once more.
\end{tcolorbox}
\section[\texorpdfstring{\kn{ಮಹಲು} (mahalu)}{ಮಹಲು (mahalu)}]{\kn{ಮಹಲು} (mahalu) — a mansion; a palace hall}

\label{lesson:KA-C371-mahalu}



\begin{quote}
Before the new one: say the Kannada for a palanquin, then the Kannada for a chieftain.
\end{quote}

\subsection*{You'll want to know: \kn{ಮಹಲು}}
\textbf{\kn{ಮಹಲು}} — \emph{mahalu} — \textquotedblleft{}a mansion; a palace hall\textquotedblright{}.

\begin{quote}
\kn{ಪಾಳೆಗಾರನ} \kn{ಮಹಲು} \kn{ಈಗ} \kn{ಒಂದು} \kn{ವಸ್ತುಸಂಗ್ರಹಾಲಯ}.
\end{quote}

\begin{quote}
\emph{pāḷegārana mahalu īga ondu vastusaṅgrahālaya.}
\end{quote}

\begin{quote}
The chieftain's mansion is now a museum.
\end{quote}

Plural: \kn{ಮಹಲುಗಳು}, \emph{mahalugaḷu}.

A large grand building; also a storey: \kn{ಎರಡು} \kn{ಮಹಲಿನ} \kn{ಮನೆ}.

\begin{cousinweb}
\textbf{\kn{ಮಹಲು}} (\emph{mahalu}) is a loan from Arabic \emph{maḥall}, \textquotedblleft{}a place, a palace, a quarter\textquotedblright{}, which came through Persian and Hindustani.
\end{cousinweb}

\begin{grammarlens}[title={What you've built}]
One more word for telling what happened.
\end{grammarlens}

\subsection*{Your turn}
\begin{itemize}
\raggedright
  \item \emph{Say it:} \emph{mahalu}
  \item \emph{Say it:} \emph{mahalu}, once more
  \item \emph{Say it:} \emph{pāḷegārana mahalu īga ondu vastusaṅgrahālaya.} — The chieftain's mansion is now a museum.
  \item \emph{Recall:} say the Kannada for a palanquin, then the Kannada for a chieftain, then say \emph{mahalu} again
\end{itemize}

\begin{tcolorbox}[breakable,colback=teal!4,colframe=teal!35!black,title={Before you move on}]
What does \kn{ಮಹಲು} mean? (\textquotedblleft{}A mansion; a palace hall\textquotedblright{}.) Say it once more.
\end{tcolorbox}
\section[\texorpdfstring{\kn{ಕಾಲಿಡು} (kāliḍu)}{ಕಾಲಿಡು (kāliḍu)}]{\kn{ಕಾಲಿಡು} (kāliḍu) — to set foot}

\label{lesson:KA-C371-kalidu}



\begin{quote}
Before the new one: say the Kannada for a chieftain, then the Kannada for a mansion.
\end{quote}

\subsection*{You'll want to know: \kn{ಕಾಲಿಡು}}
\textbf{\kn{ಕಾಲಿಡು}} — \emph{kāliḍu} — \textquotedblleft{}to set foot\textquotedblright{}.

\begin{quote}
\kn{ಹತ್ತು} \kn{ವರ್ಷದ} \kn{ನಂತರ} \kn{ಅವನು} \kn{ತನ್ನ} \kn{ಊರಿಗೆ} \kn{ಕಾಲಿಟ್ಟನು}.
\end{quote}

\begin{quote}
\emph{hattu varṣada nantara avanu tanna ūrige kāliṭṭanu.}
\end{quote}

\begin{quote}
After ten years he set foot in his home town.
\end{quote}

Past (I …): \kn{ಕಾಲಿಟ್ಟೆ}, \emph{kāliṭṭe}.

Where one steps takes -\kn{ಗೆ}. It also means to enter a new stage: \kn{ಹೊಸ} \kn{ವರ್ಷಕ್ಕೆ} \kn{ಕಾಲಿಡು}.

\begin{cousinweb}
\textbf{\kn{ಕಾಲಿಡು}} (\emph{kāliḍu}) is a native Dravidian compound: \emph{kāl}, \textquotedblleft{}foot, leg\textquotedblright{} (DEDR 1479), and \emph{iḍu}, \textquotedblleft{}to put, to place\textquotedblright{} (DEDR 442).
\end{cousinweb}

\begin{grammarlens}[title={What you've built}]
One more word for telling what happened.
\end{grammarlens}

\subsection*{Your turn}
\begin{itemize}
\raggedright
  \item \emph{Say it:} \emph{kāliḍu}
  \item \emph{Say it:} \emph{kāliḍu}, once more
  \item \emph{Say it:} \emph{hattu varṣada nantara avanu tanna ūrige kāliṭṭanu.} — After ten years he set foot in his home town.
  \item \emph{Recall:} say the Kannada for a chieftain, then the Kannada for a mansion, then say \emph{kāliḍu} again
\end{itemize}

\begin{tcolorbox}[breakable,colback=teal!4,colframe=teal!35!black,title={Before you move on}]
What does \kn{ಕಾಲಿಡು} mean? (\textquotedblleft{}To set foot\textquotedblright{}.) Say it once more.
\end{tcolorbox}
